\bookchapter{附錄A：新手機開箱「保姆級」流程}

\booksection{開箱驗機「三部曲」：}

\begin{itemize}
  \item \textbf{睇外觀}：檢查機身、螢幕有無劃痕、磕碰。
  \item \textbf{對資訊}：核對盒子上嘅IMEI/序列號同發票上嘅係咪一致。
  \item \textbf{試功能}：撥個電話試下聽筒話筒，拍張相檢查攝像頭。
\end{itemize}

\booksection{激活與初始設定「重點提醒」：}

\begin{itemize}
  \item \textbf{網絡}：先連上穩定嘅Wi-Fi。
  \item \textbf{帳號}：登入你嘅蘋果ID、華為帳號、小米帳號等。
  \item \textbf{基礎設定}：將「查找我嘅手機」打開，設定一個靠譜嘅鎖屏密碼。
  \item \textbf{eSIM設定}：如果你嘅手機支持eSIM（例如iPhone、部分安卓旗艦），可以喺設定入面添加eSIM套餐，無需插入實體SIM卡。出國嗰陣eSIM尤其方便，可以直接購買當地運營商嘅eSIM套餐。
  \item \textbf{數據遷移}：由舊手機遷移數據嗰陣，確保兩台手機電量充足。iPhone用「快速開始」，安卓用官方換機App。
\end{itemize}

\booksection{裝App「安全第一」：}

\begin{itemize}
  \item \textbf{首選官方商店}：蘋果用App Store，安卓用手機自帶嘅應用商店。
  \item \textbf{權限唔好亂畀}：安裝嗰陣彈出嘅權限請求，諗清楚「呢個App真係需要嗎？」
\end{itemize}
